% https://github.com/MCG-NKU/NSFC-LaTex
% by Ming-Ming Cheng https://mmcheng.net
% 关于VsCode LaTeX的配置 https://www.cnblogs.com/ourweiguan/p/11785660.html
% Windows下Tex Live最新版可以用pdfLatex快速编译和标准模板相似度更高的文档。



\documentclass[12pt]{article}


\usepackage{nsfc}
\setmainfont{Times New Roman}

%\usepackage{fontspec}
%\usepackage{xcolor}
%\defaultfontfeatures{Ligatures=TeX}




\newcommand{\addImg}[2][1.0]{\includegraphics[width=#1\linewidth]{#2}}
\newcommand{\cmm}[1]{\textcolor[rgb]{0,0.6,0}{CMM: #1}}
\newcommand{\todo}[1]{{\textcolor{red}{\bf [#1]}}}
\newcommand{\myPara}[1]{\paragraph{#1：}}
\newcommand{\myEmph}[1]{\textbf{\textcolor[rgb]{0,0,0.25}{#1}}}
\newcommand{\mySec}[1]{\vspace{0.15in} \noindent \myEmph{$\square$~#1}}

\graphicspath{{figures/}}


\begin{document}




%\if 0
\begin{center}\bf\Large 开放环境下3D点云分析测试时自适应方法研究 \\

Research on Test-time Adaptation for 3D Point Cloud Analysis in Open World 
\end{center}





\textbf{摘要}：
%
“用……方法（手段）进行……研究，探索/证明……问题，
对阐明……机制／揭示……规律有重要意义，为……奠定基础／提供……思路 ”。
作用：画龙点睛。
效果：引发评议专家兴趣，使其产生探个究竟的好奇心——“他到底要怎么做？”
摘要字少，但切忌平淡无奇（要勾起评委浓厚兴趣）。
\myEmph{一定要语气坚定，旗帜鲜明，字数有限，资源宝贵，要特别注意重点突出，
讲明现状、意义、课题主要研究目标、内容、思路和预期结果}。



\textbf{Abstract}：
***


科学问题属性：聚焦前沿，独辟蹊径



【待补充】


\clearpage
%\fi


%%%%%%%%% TITLE
\title{报告正文}

\maketitle
\thispagestyle{empty}

%\emph{\large 参照以下提纲撰写，要求内容翔实、清晰，层次分明，标题突出。
%\nsfClr{请勿删除或改动下述提纲标题及括号中的文字。}}


\ContentDes{（一）立项依据与研究内容\textnormal{\alert{（建议$8000$字以内）}}：}


\NsfcSection{1}{项目的立项依据}{
（研究意义、国内外研究现状及发展动态分析，需结合科学研究发展趋势来论述科学意义；
或结合国民经济和社会发展中迫切需要解决的关键科技问题来论述其应用前景。
附主要参考文献目录）；}

% 写作思路是什么
% 1. 倒推 
% 1.1 本研究到底要解决什么问题？    ->  开放环境下点云处理方法效果差的问题，这里最好用数据说明
% 1.2 我提出了什么方案？           ->  测试时自适应缓存模型
% 1.3 提出的方案有何好处？         ->  拔插即用/免训练/高效/在线性能提升

% 2. 正推
% 2.1 简述研究背景
% 2.2 该背景下重点关注什么任务？
% 2.3 该任务存在什么问题？
% 2.4 本项目的解决思路      ->      主要讲思想
% 2.5 本项目的技术方案      ->      主要讲技术路径
% 2.6 提出的解决方案有何优势？
% 2.7 提出的解决方案有何科学意义和重要价值


\subsection{研究背景及科学意义}

% 从这几方面展开，提供参考文献
% 讲开放环境点云处理
%   讲明白过去的范式是怎么样的？特定数据集上训练模型，特定数据集上测试
%   讲明白现在朝着开放世界模型走？
%   存在的问题：鲁棒性/泛化性差/效率低
2017年，以PointNet~\cite{qi17pointnet,qi17pointnet2}为代表的深度点云处理方法横空出世，它们在点云物体识别、语义分割等
任务上的性能大幅超越人工特征工程和传统机器学习方法，（用ModelNet举例子，比传统方法好多少）。经过多年的发展，深度点云处理方法在
特定数据集（封闭设定）的性能趋于饱和（给出具体例子，发表了多少篇论文，性能增益不到 6\%），
从2022年开始，越来越多的工作关注开放环境下点云处理问题，比如开放词表的点云识别~\cite{zhang22pointclip,zhu23pointclip2,xue23ulip,xue24ulip2,liu23openshape,zhou24uni3d}、
物体检测~\cite{zeng23clip2,zhu23pointclip2,lu23ov3det,cao23coda,yang24imov3d,cao25codav2}、语义分割，
当前该领域的研究重点已经从封闭数据集的设定迈向开放环境，且逐渐成为国内外一个前沿研究问题。

% 进一步讲开放环境的点云理解现状，存在什么问题
%   存在的挑战一句话就说完了，得扩充和丰富
但当前开放环境下点云处理方法未能良好应对千变万化的数据，模型的准确性和鲁棒性非常脆弱（这里要举例子，至少目前看到这么多点云识别/检测的例子，准确性只有那么点），
不能满足实际应用需求。例如，当前最好的开放词表的点云物体检测模型 CoDA v2~\cite{cao25codav2}，在公开的ScanNet数据集上的平均精度 mean Average Precision (mAP@IoU$_{25}$) 仅为 22.72\%，
而在封闭数据集上训练的3DETR~\cite{misra213detr}模型在ScanNet的 mAP@IoU$_{25}$ 达到62.7\%，当前经过封闭数据训练的最佳方法Dest~\cite{wang25dest}在ScanNet达到了 78.3\% mAP@IoU$_{25}$ 的检测精度，领先CoDA v2 55.58\%个绝对点。

% 讲本项目思路
%   解释本项目涉及到的关键技术，就是下面列的
针对这些问题，本项目提出开放环境下基于缓存模型的测试时自适应点云分析方法，致力于提升开放环境下点云分析方法的鲁棒性和泛化能力。

% 讲测试时自适应
%   讲主要思想，这要结合文献
尽管模型在训练阶段要尽可能学习各种情况，但现实世界是千变万化的，总有些新的数据特性在测试和部署阶段出现，
此时训练好的模型不再具备大规模重训练可行性，对于测试环境的应变能力和自适应能力显得尤为重要。
而测试时自适应技术旨在让模型在线运行时具备适应测试环境或根据测试数据特点调整自身预测的能力~\cite{yeo23rna}。
%   讲技术分类，也要总结文献
测试时自适应技术可大致分为两类，一类是需要根据测试样本调整预训练模型参数的~\cite{schneider20improving,wang21tent,zhang21adaptive,zhang22memo,shu22tpt,feng23difftpt,niu23towards,chen22adacontrastive,zhao2024rlcf}，另外一类是免训练的~\cite{wang24bftt3d,karmanov24tda,zanella24meanshift,farina24zero,sun25pointcache,tamjidi2025uniadapter}。在第一大类方法中，有调整BN层统计信息的~\cite{schneider20improving,wang21tent,zhang22memo}，有调整提示词的~\cite{shu22tpt,feng23difftpt}，还有调整模型其他部分参数（如图像编码器）的~\cite{zhang21adaptive,chen22adacontrastive,zhao2024rlcf}；在第二大类方法中，大部分是基于缓存模型的~\cite{wang24bftt3d,karmanov24tda,sun25pointcache,tamjidi2025uniadapter}，还有些经过理论推导进行数值变换的~\cite{zanella24meanshift,farina24zero}。

% 讲缓存模型——测试时自适应方法的一种
%   讲主要思想，就这一句话吗？不止这样吧，没把精华讲出来
缓存模型通过存储典型样本数据为模型在评测和部署时作参考。
%   讲技术路线分类？
这种思路和技术手段已在多种任务重得到体现和应用，比如语言建模~\cite{vinyals16mn,grave17unbounded,merity17pointer,khandelwal20generalization}、图像识别~\cite{vinyals16mn,snell17prototypical,finn17maml,chen18a,chen21meta-baseline,choi22improving,iwasawa21t3a,zhang22tip-adapter,zhang23cafo,karmanov24tda}、点云处理~\cite{zhang23pointnn,tang24pointpeft,sun25pointcache,wang24bftt3d,tamjidi2025uniadapter}。
%   讲本文如何看待缓存模型  ->  检索增强范式

% 讲存在的问题？
%   评述总结，该把这部分写出来了，不能磨了
%   1. TTA 存在很多假设，限制了实际应用
%   2. TTA 效率低，不能实时运行
%   3. 缓存模型依赖训练数据构建
然而，当前的测试时自适应点云分析方法面临以下几方面挑战：
%   分别对应后面的主要研究内容
%   最后一定要有对挑战的总结描述
\begin{itemize}
    \item 当前测试时自适应点云分析方法大多限定在封闭数据集设置，仅能处理训练过程已经见过的物体和场景类型，不具备应对开放环境的自适应能力。
    \item 当前测试时自适应点云分析方法依赖大规模训练数据构建缓存，但是模型测试和部署阶段获取大量训练数据不太实际，而且测试数据展示的特性和训练数据可能很不相同，导致这类方法的鲁棒性和泛化能力较弱，且运行效率不高，难以投入实际应用。
    \item 当前开放环境下测试时自适应点云分析方法的研究处于早期阶段，基础理论尚不完备，关键方法和技术路线尚不成熟。
\end{itemize}

综上所述，开放环境下测试时点云处理的研究处于起步阶段，相关基础理论有待发展，对应技术路径亟待验证。本项目计划从理论和实践两个维度展开研究：

（1）从理论角度出发，开放环境下测试时自适应点云分析方法理论框架尚不清晰。

（2）从应用角度出发，缓存模型构建和更新不够准确和高效。

（3）从应用角度出发，开放环境下基于缓存模型的测试时自适应点云分析方法未形成体系。

\subsection{本人具备的研究基础}

本人博士期间围绕3D点云处理的准确性、运行效率、鲁棒性和泛化性进行深入研究，广泛阅读了领域文献，熟悉并掌握了主要技术路径和方法，提出了系列创新方案，以第一作者身份围绕该主题发表了5篇顶级国际会议和期刊论文~\cite{sun25pointcache,sun24point-prc,sun25vsformer,sun24ppt3d,sun23vipformer}，初步具备了执行本项目的研究基础和实践经验。
其中，发表在CVPR 2025的论文(Point-Cache)探讨了开放环境下点云识别模型测试时鲁棒性和泛化性问题；发表在 NeurIPS 2025 的论文(Point-PRC)探讨了3D大模型的领域泛化能力；发表在 ICRA 2024 的论文(PPT) 探讨了如何用参数和数据高效的方式激活多模态3D大模型的预训练知识，从而增强点云处理的效果和效率；发表在 ICRA 2023 的论文探讨了如何结合图像和点云两个模态的信息增强点云处理的性能和效率。

% 研究基础总结
本人将充分利用这些工作积累，深入研究开放环境下3D点云模型测试时自适应问题，努力突破其中的关键科学问题和技术挑战，给出适合3D点云大模型自适应部署的理论和实践方案，为Spatial AI的发展贡献力量。本项目预期产出的研究成果对于3D场景重建、3D物体生成、3D物体配准、VR/AR等领域也具有重要参考价值。

\subsection{国内外相关工作}

\textbf{开放环境下3D点云处理}
% 从哪讲起呢？把事情发展的原委讲出来，配上自身评述

近几年，以CLIP~\cite{radford21clip}和ALIGN为~\cite{jia21scaling}代表的视觉-语言对比预训练方法受到了广泛关注，在这种思路指导下训练的模型能进行开放集合的图像理解（如图像识别/语义分割/物体检测），并且表现出显著改善的泛化能力，极大扩展了之前仅在图像单一模态训练的模型的能力边界（后者仅能处理训练阶段见过的物体/语义类别）。受到这类工作的启发，开放环境下3D点云处理的研究逐渐兴起，当前呈现出蓬勃发展的趋势。

% 怎么和上述概览的逻辑接起来
%   介绍典型任务分类
开放环境下3D点云处理典型的任务包括3D点云识别~\cite{zhang22pointclip,zhu23pointclip2,sun25pointcache,sun24point-prc,sun24ppt3d,xue23ulip,xue24ulip2,liu23openshape,zhou24uni3d}、
% PointCLIP~\cite{zhang22pointclip},Point-Cache~\cite{sun25pointcache}, Point-PRC~\cite{sun24point-prc}, PPT~\cite{sun24ppt3d}, ULIP~\cite{xue23ulip}, ULIP-2~\cite{xue24ulip2}, OpenShape~\cite{liu23openshape}, Uni3D~\cite{zhou24uni3d}
3D点云物体检测~\cite{zeng23clip2,zhu23pointclip2,lu23ov3det,yang24imov3d,cao23coda,cao25codav2}、
% CLIP2~\cite{zeng23clip2} (Clip2: Contrastive language-image point pretraining from real-world point cloud data), PointCLIP V2~\cite{zhu23pointclip2}, OV-3DET~\cite{lu23ov3det}, ImOV3D~\cite{yang24imov3d}, CoDA~\cite{cao23coda}, CoDA v2~\cite{cao25codav2}
3D点云语义分割~\cite{jiang24ov3d,kolodiazhnyi2024oneformer3d,zhou2025pointsam,rusnak25haec,boudjoghra2025openyolo3d}等。
% OV3D: Open-Vocabulary 3D Semantic Segmentation with Foundation Models~\cite{jiang24ov3d}
% OneFormer3D~\cite{kolodiazhnyi2024oneformer3d}, 
% Point-SAM~\cite{zhou2025pointsam}, 
% HAECcity~\cite{rusnak25haec}
% Open YOLO 3D~\cite{boudjoghra2025openyolo3d}, 
%   总结每类任务特点 有什么特点呢？这就要自己琢磨了
这几类任务起来有很大差异，从数据规模和复杂度层面，3D点云识别涉及到物体级的点云数据，规模和复杂度相对较小，而3D物体检测和语义分割要处理场景级的点云数据，规模和复杂度较大，因此不同任务的处理方法也有较大差异。

%   评述这些工作与本项目的联系和区别
然而，上述开放集合的点云处理方法缺乏自适应能力，无法根据测试环境特点作出自适应调整，极大限制了3D点云处理方法在开放环境和真实世界的部署应用。本项目正是针对这种缺陷，研究开放环境下3D点云处理的自适应方法，让3D模型能够高效的根据测试数据特点调整自身预测，从而改善模型的准确性和鲁棒性。

\textbf{测试时自适应方法}
% 根据Point-Cache这篇文章的思路，分类总结
%   分成哪几类呢？
%   - 2D TTA：
%       1. 非开放环境   <-  基于单模态   
%       2. 开放环境     <-  多模态
%       - 用不用按照实现方案把技术分个类？上面的2类太简单？
%   - 3D TTA
%       1. 基于小模型：鲁棒性和泛化能力本身不强+不是开放环境
%       2. 依赖庞大训练集构建pipeline
%       3. 未利用测试数据分布

测试时自适应技术在图像处理领域已经得到体现和应用。早期的图像处理测试时自适应技术应用在图像单一模态~\cite{liang20shot,liang22source,wang21tent,zhang22memo,sun20ttt,liu21ttt++,gandelsman22ttt-mae,iwasawa21t3a,zhang22memo}，由于未引入灵活的语言模态，这类方法只能在封闭数据集上进行图像识别，不具备在真实开放世界中进行图像处理的能力。近年来，测试时自适应技术与主流视觉-语言模型结合，掀起了开放环境下多模态大模型测试时自适应方法研究的浪潮~\cite{shu22tpt,feng23difftpt,tsai24gda,karmanov24tda,zanella24meanshift,farina24zero,guan2025sca,cha2025snap}，这些工作大大突破了封闭数据集合的设定，使得视觉-语言多模态模型更加贴近实际应用。

% ====== 2D TTA: uni-modal, close set ======
    % Memo~\cite{zhang22memo}

    % Test-time classifier adjustment module for model-agnostic domain generalization~\cite{iwasawa21t3a} (neurips 2021)

    % Tent: Fully test-time adaptation by entropy minimization~\cite{wang21tent} (iclr 2020),

    % Do we really need to access the source data? Source hypothesis transfer for unsupervised domain adaptation.~\cite{liang20shot} (icml 2020),

    % Source data-absent unsupervised domain adaptation through hypothesis transfer and labeling transfer.~\cite{liang22source} (TPAMI 2022),

    % TTT: Test-time training with self-supervision for generalization under distribution shifts.~\cite{sun20ttt}

    % TTT++~\cite{liu21ttt++} (neurips 2021),
    % The norm must go on: Dynamic unsupervised domain adaptation by normalization.~\cite{mirza22dua} (CVPR 2022),

    % TTT-MAE~\cite{gandelsman22ttt-mae} (Test-Time Training with Masked Autoencoders, neurips 2022) 

% ====== 2D TTA: multi-modal, open set ======
    % TPT, DiffTPT, TDA, 
    % GDA: Generalized diffusion for robust test-time adaptation.~\cite{tsai24gda} (CVPR 2024), 

    % On the test-time zeroshot generalization of vision-language models: Do we really need prompt learning?~\cite{zanella24meanshift} (CVPR 2024)

    % Zero~\cite{farina24zero}, SCA~\cite{guan2025sca}, SNAP~\cite{cha2025snap}

与在图像处理领域的探索相比，3D点云测试时自适应技术起步较晚，相关研究仍属于早期阶段。当前的技术流派可大致划分为三大类：第一类是基于小模型的方法~\cite{mirza23mate,wang24bftt3d,shim24cloudfixer}，这类方法首先仅局限于点云模态，缺乏开放集合的3D点云处理能力，其次由于模型本身能力有限，虽有测试时自适应技术加持，但展现的准确性和泛化能力不够；第二类是基于多模态大模型的方法~\cite{sun25pointcache,tamjidi2025uniadapter}，这类方法虽然能进行开放集合的点云处理，但是没有充分利用测试数据的分布特性（前/后面这段话是 **placehoder**），导致模型在现实世界测试和部署时鲁棒性和泛化能力仍不理想。

% ====== 3D TTA ======
% MATE, BFTT3D, CloudFixer, Point-Cache, Uni-Adapter

% 相关工作总结，本项目的特色之处

\textbf{缓存模型(检索增强方法，这么说不准确)}
% 这部分放到测试时自适应中讲？还是单拎出来？
% 目前标题中已经没有“缓存模型”的字样了，单拎出来有点突兀

语言建模方面：
Unbounded cache model for online language modeling with open vocabulary, neurips 2017

Generalization through memorization: Nearest neighbor language models, iclr 2020

Pointer sentinel mixture models, iclr 2017

图像识别方面：
Matching networks for one shot learning, neurips 2016

A closer look at few-shot classification, iclr 2019

Meta-baseline: Exploring simple meta-learning for few-shot learning, ICCV 2021

Improving test-time adaptation via shift-agnostic weight regularization and nearest source prototypes, eccv 2022

A simple cache model for image recognition. neurips 2018

Prototypical networks for few-shot learning. neurips 2017

Sus-x: Training-free name-only transfer of vision-language models, ICCV 2023

Tip-Adapter, eccv 2022 CaFo, CVPR 2023

Test-time classifier adjustment module for model-agnostic domain generalization, neurips 2021
TDA, CVPR 2024

3D方面：PointNN, Point-PEFT, BFTT3D, Point-Cache, Uni-Adapter

\myEmph{采纳这个建议}

1. 可以引用几个代表性图，形象的展示一些重要的背景知识，方便大同行理解。

2. 可以对相关工作进行分组介绍。

3. 每组介绍之后简单总结一下现有工作和拟研究工作的关系。

\myEmph{现有工作存在那些不足需要进一步研究。
这些总结性的结论建议粗体强调，方便评审人迅速理解。}

{
\bibliographystyle{nsfc.bst}
\bibliography{nsfc2026}
}


%%%%%%%%%%%%%%%%%%%%%%%%%%%%%%%%%%%%%%%%%%%%%%%%%
\NsfcSection{2}{项目的研究内容、研究目标，以及拟解决的关键科学问题}{
（此部分为重点阐述内容）；}

\subsection{研究目标}


重点阐述本项目计划\myEmph{做到什么程度}。
一定要用大同行能够理解的术语来描述。


\subsection{研究内容}


重点阐述本项目计划\myEmph{做什么}。
如\figref{fig:teaser}所示，建议撰写之前仔细画一个主要研究内容的图。
我通常习惯用PowerPoint作图，做好后导出为pdf格式，既可以保持图片不会文件太大，也可以保证放大后非常清晰。
幻灯片直接导出的pdf可能存在空白边，可以用WPS中的“页面-剪裁页面”去掉白边。
作图用的pptx文件我也通常会保存起来，例如这个模版\LaTeX 文件中的“prepare/NSFC-Figs.pptx”。
如果是大的国基金项目，后续可能涉及答辩。
答辩时这个图可以用，保留pptx格式也方便到时候做尺寸和布局的调整。


\begin{figure}[ht]
	\centering
    \begin{overpic}[width=0.8\columnwidth]{framework.pdf}
    \end{overpic}
    \caption{本项目主要研究内容。
    }\label{fig:teaser}
\end{figure}




\subsection{拟解决的关键科学问题}



\NsfcSection{3}{拟采取的研究方案及可行性分析}{
（包括研究方法、技术路线、实验手段、关键技术等说明）；}


\subsection{拟采取的技术路线}

如图\figref{fig:pipline}所示，建议在具体写技术路线之前，先厘清这个框架图。
技术路线的框架图通常比研究内容更加饱满，可以更好的展示本项目的研究思路。

\begin{figure}[ht]
	\centering
    \begin{overpic}[width=\columnwidth]{pipeline.pdf}
    \end{overpic}
    \caption{本项目的技术路线。
    }\label{fig:pipline}
\end{figure}


\subsection{可行性分析}

既然国内外相关工作都没能解决你提出的重要问题，为什么你觉得自己有望解决该问题。
论述的时候通常包括：独特的时机、与众不同的方案、雄厚的相关科研积累等。


\NsfcSection{4}{本项目的特色与创新之处；}{}

您的方法有什么独特性，为什么您认为它会成功？

\NsfcSection{5}{年度研究计划及预期研究结果}{
（包括拟组织的重要学术交流活动、国际合作与交流计划等）。}

\subsection{年度研究计划}


\subsection{预期研究成果}

项目各阶段，特别是中期和期末，如何检查该项目计划成功与否？

%%%%%%%%%%%%%%%%%%%%%%%%%%%%%%%%%%%%%%%%%%%%%%%%%
\ContentDes{（二）研究基础与工作条件}


\NsfcSection{1}{研究基础}{
（与本项目相关的研究工作积累和已取得的研究工作成绩）；}

申请人自2017年开始从事基于深度学习的计算机视觉和三维视觉相关研究工作，具备扎实的理论基础和丰富的实践经验，取得了一些理论和应用创新成果。申请人的博士学位论文《多视图多模态三维物体理解关键技术研究》获评中国人民大学2025年度优秀博士学位论文，申请人于2023年获得国家留学基金委公派留学奖学金,后赴澳大利亚莫纳什大学（Monash University）数据科学与人工智能系进行为期一年的博士联合培养，联培期间的研究方向仍是三维视觉和点云处理，合作导师为蔡剑飞教授（Prof. Jianfei Cai），与蔡剑飞教授团队建立了良好而深入的合作关系。

本人博士期间围绕3D点云处理的准确性、运行效率、鲁棒性和泛化性进行深入研究，广泛阅读了相关文献，熟悉主要技术路径和方法，提出系列创新方案并发表了系列高水平论文，在该领域具备扎实的研究基础。例如，本人以第一作者发表在 CVPR 2025 的论文提出《Point-Cache: Test-time Dynamic and Hierarchical Cache for Robust and Generalizable Point Cloud Analysis》提出了一个动态和分层缓存模型，大幅改善了3D点云处理模型在测试/部署阶段的鲁棒性和泛化能力；以第一作者发表在 NeurIPS 2024 的论文《Point-PRC: A Prompt Learning Based Regulation Framework for Generalizable Point Cloud Analysis》提出了一种基于提示学习的即插即用的监管框架，该框架能无缝集成到不同多模态3D大模型中，在增强点云处理下游任务性能的同时，改善多模态3D大模型在未知点云数据的泛化能力；以第一作者发表在 TVCG 2025 的论文《VSFormer: Mining Correlations in Flexible View Set for Multi-view 3D Shape Understanding》设计了一种灵活的3D物体多视图关系建模方法，增强了不同视图间的信息流动和交互，提升了3D物体建模和识别的准确性；以第一作者发表在 ICRA 2024 的论文《Parameter-efficient Prompt Learning for 3D Point Cloud Understanding》提出了一种参数和数据双重高效的提示学习方法，仅用极少训练参数和数据样本就能激活多模态3D大模型丰富的预训练知识，显著提升了模型在点云处理任务中的准确性和运行效率；以第一作者发表在 ICRA 2023 的论文《ViPFormer: Efficient Vision-and-Pointcloud Transformer for Unsupervised Pointcloud Understanding》提出了一种无监督预训练模型架构，该架构能有效融合图像和点云多模态数据信息，促进点云处理下游任务取得更好效果，同时提升下游任务的运行效率。

以我本人的第一个国基金项目的\myEmph{申请原文为例}，通过高度精炼的语言，
在2页左右的篇幅中快速讲清楚：1）大部分申请人没有的特色优势是什么；
2）特色工作和拟研究工作的联系得清晰。

申请人在与本项目相关的研究工作中有着丰富的积累，并取得了国际领先的科研成果, 
发表\myEmph{10余篇CCF A类}国际期刊及会议论文（ACM TOG 4篇，IEEE TPAMI 2篇，
IEEE CVPR 3篇，IEEE ICCV 2篇，IEEE TVCG 2篇），
\myEmph{论文Google Scholar他引1600+次，一作论文单篇最高他引700+次}。
这些研究基础主要包括3个方面：


\subsection{图像的智能理解与新一代交互方面}


\subsection{显著性物体检测方面}


\subsection{基于视觉注意机制的候选物体生成方面}



\subsection{研究工作获奖}


\begin{figure}[ht]
    \centering
    \addImg[.8]{figures/awards.jpg}
    \caption{图像中候选物体（object proposal）生成的快速机制。}
    \label{fig:award}
\end{figure}

申请人在图像分析与编辑方面的初步工作获得了多项重要奖励，
包括：Google PhD Fellowship (2010年亚洲地区共2人获奖)，
北京市优秀博士论文（同年度清华大学56个系共有7人获得该奖项），
英国皇家学会Newton International Fellowship提名，教育部自然科学奖一等奖，
IBM PhD fellowship，北京市优秀博士毕业生，教育部博士研究生学术新人奖等奖项。
这些初步研究中取得的成果为本项目取得进一步的创新性成果打下了坚实的基础。


\NsfcSection{2}{工作条件}{
（包括已具备的实验条件，尚缺少的实验条件和拟解决的途径，
包括利用国家实验室、
全国重点实验室和部门重点实验室等研究基地的计划与落实情况）；}


\myPara{申请平台概况}
南开大学计算机学科在计算机视觉与计算机图形学领域具有非常扎实的研究基础。
作为教育部直属重点大学，南开大学是国内学科门类最齐全的综合性、研究型大学之一。
在计算机视觉与图形学方向上，
南开大学计算机学院拥有计算机与控制工程国家级虚拟仿真实验教学中心、
可信行为智能算法与系统教育部工程研究中心、和天津市视觉计算与智能感知重点实验室
等一系列优良研究平台。
计算机视觉与图形学团队现拥有一批先进的高性能GPU服务器集群（包含Tesla A40, Tesla V100, RTX 3090等高端GPU共计900多块）。
% 
为本项目的研发提供强有力的计算环境。
南开大学为计算机视觉与图形学团队提供了1200平米的科研用房。
这些良好的配套支持将为项目的顺利开展提供优良的工作条件。


\myPara{科研团队介绍}

申请人所在的南开大学计算机视觉与图形学团队
由国家杰出青年基金获得者程明明教授带领，
包含国家级“四青”人才5人。
%
近五年，团队承担国家自然科学基金重点项目、国家重点研发计划课题、
国防科技创新重点项目等重点重大项目10余项；
在相关领域的 SCI 一区/CCF A 类顶级国际期刊和会议上发表学术论文100余篇，
其中 TPAMI论文 40余篇，ESI高被引论文30余篇；
获得教育部自然科学一等奖、中国图象图形学学会自然科学一等奖、
吴文俊人工智能自然科学二等奖等多项奖励。


\myPara{国内外合作与学术交流情况}

南开大学计算机视觉与图形学团队有着丰富的国内外学术合作与交流基础。
近年来，团队与计算机视觉最高奖Marr奖得主、英国皇家学会院士、
牛津大学Philip Torr教授团队，
计算机视觉最高奖Marr奖得主、加州大学圣迭戈分校的Zhuowen Tu教授团队，
和清华大学胡事民教授团队合作开展了多项具有影响力的学术研究，
并共同发表相关学术论文。
团队成员受邀担任IEEE TPAMI, IEEE TIP, IEEE CVPR, ICCV等本领域
多个顶级期刊和会议的编委或领域主席。
团队成员作为主要组织者承办了VALSE 2022 (3000余人现场参会)，2023 (5000余人现场参会)，
和PRCV 2020，2024 (1000多人现场参会)
等多个国内大型学术会议。
%
与这些国内外相关研究团队开展密切合作与交流的经验也将为项目的成功
进行起到重要的促进作用。

申请人所申报的项目将依托于南开大学计算机视觉与图形学团队开展。
现有优越的实验条件可以充分满足项目研发需求。
项目团队将继续加强国际交流与合作，
以创新性的高水平学术成果为研究目标。


\NsfcSection{3}{正在承担的与本项目相关的科研项目情况}{
（申请人正在承担的与本项目相关的科研项目情况，包括国家自然科学基金
的项目和国家其他科技计划项目，要注明项目的资助机构、项目类别、
批准号、项目名称、获资助金额、起止年月、与本项目的关系及负责
的内容等）；}

正在承担的项目

\NsfcSection{4}{完成国家自然科学基金项目情况}{
（对申请人负责的前一个已资助期满的科学基金项目（项目名称及批准号）完成情况、后续研究进
展及与本申请项目的关系加以详细说明。另附该项目的研究工作总结
摘要（限500字）和相关成果详细目录）。}

已完成的国家自然科学基金

%%%%%%%%%%%%%%%%%%%%%%%%%%%%%%%%%%%%%%%%%%%%%%%%%
\ContentDes{（三） 其他需要说明的情况}



\NsfcSection{1}{}{申请人同年申请不同类型的国家自然科学基金项目情况（列明
同年申请的其他项目的项目类型、项目名称信息，并说明与本项目之
间的区别与联系；已收到自然科学基金委不予受理或不予资助决定的，
无需列出）。}


\NsfcSection{2}{}{具有高级专业技术职务（职称）的申请人是否
存在同年申请或者参与申请国家自然科学基金项目的单位不一致的情
况；如存在上述情况，列明所涉及人员的姓名，申请或参与申请的其
他项目的项目类型、项目名称、单位名称、上述人员在该项目中是申
请人还是参与者，并说明单位不一致原因。}



\NsfcSection{3}{}{具有高级专业技术职务（职称）的申请人是否
存在与正在承担的国家自然科学基金项目的单位不一致的情况；如存
在上述情况，列明所涉及人员的姓名，正在承担项目的批准号、项目
类型、项目名称、单位名称、起止年月，并说明单位不一致原因。}


\NsfcSection{4}{}{同年以不同专业技术职务（职称）申请或参与申请科学基金项
目的情况（应详细说明原因）。}

\NsfcSection{5}{}{其他。}

无


\end{document}
